\documentclass[12pt]{article}

\usepackage{fullpage}
\usepackage{graphicx}
\usepackage{amssymb}
\usepackage{amsmath}
\usepackage[none]{hyphenat}
\usepackage{parskip}
\usepackage[spanish]{babel}
\usepackage[utf8]{inputenc}
\usepackage{hyperref}
\usepackage{fancyhdr}
\usepackage{tasks}
\usepackage{mdframed}
\usepackage{xcolor}
\usepackage{pgfplots}
\usepackage[makeroom]{cancel}
\usepackage{multicol}
\usepackage[shortlabels]{enumitem}
\usepackage{stackrel}
\usepackage{tkz-tab}
\usepackage{xpatch}
\xpatchcmd{\tkzTabLine}{$0$}{$\bullet$}{}{}

\setlength{\headheight}{10pt}
\setlength{\headsep}{10pt}
\pagestyle{fancy}
\rhead{\ayudantia \ - \alumno}
\tikzset{t style/.style={style=solid}}

\newcommand*{\mybox}[2]{\colorbox{#1!30}{\parbox{.98\linewidth}{#2}}}

\newenvironment{solucion}
{\begin{mdframed}[backgroundcolor=black!10]
		{\bf Solución:}\\
	}
	{
	\end{mdframed}
}

\newenvironment{alternativas}[1]
{\begin{multicols}{#1}
		\begin{enumerate}[a)]
		}
		{
		\end{enumerate}
	\end{multicols}
}

\newenvironment{preguntas}
{\begin{enumerate}\itemsep12pt
	}
	{
	\end{enumerate}
}

\newcommand{\ayudantia}{{\sc Ayudantía 2}}
\newcommand{\tituloayu}{Integrales inmediatas y sustitucion}
\newcommand{\fecha}{19 de marzo de 2018}
\newcommand{\sigla}{MAT220E}
\newcommand{\nombre}{Cálculo II}
\newcommand{\profesor}{Vania Ramirez}
\newcommand{\ano}{2018}
\newcommand{\semestre}{1}
\newcommand{\mail}{mat220e@ifcastaneda.cl}
\newcommand{\alumno}{Ignacio Castañeda - \mail}

\newcommand{\ev}{\Big|}
\newcommand{\ra}{\rightarrow}
\newcommand{\lra}{\leftrightarrow}
\newcommand{\N}{\mathbb{N}}
\newcommand{\R}{\mathbb{R}}
\newcommand{\Exp}[1]{\mathcal{E}_{#1}}
\newcommand{\List}[1]{\mathcal{L}_{#1}}
\newcommand{\EN}{\Exp{\N}}
\newcommand{\LN}{\List{\N}}
\newcommand{\comment}[1]{}
\newcommand{\lb}{\\~\\}
\newcommand{\eop}{_{\square}}
\newcommand{\hsig}{\hat{\sigma}}
\newcommand{\widesim}[2][1.5]{
	\mathrel{\overset{#2}{\scalebox{#1}[1]{$\sim$}}}
}
\newcommand{\wsim}{\widesim{}}
\newcommand{\lh}{\stackrel{L'H}{=}}

\begin{document}
\thispagestyle{empty}

\begin{minipage}{2cm}
	\includegraphics[width=2cm]{../../../../img/logo.pdf}
	\vspace{0.5cm}
\end{minipage}
\begin{minipage}{\linewidth}
	\begin{tabular}{lrl}
		{\scriptsize\sc Pontificia Universidad Catolica de Chile} & \hspace*{0.7in}Curso: &
		\sigla  - \nombre\\
		{\sc Facultad de Matemáticas}&
		Profesor: & \profesor \\
		{\sc Semestre \ano-\semestre} & Ayudante: & {Ignacio Castañeda}\\
		& {Mail:} & \texttt{\mail}
	\end{tabular}
\end{minipage}

\vspace{-10mm}
\begin{center}
	{\LARGE\bf \ayudantia}\\
	\vspace{0.1cm}
	{\tituloayu}\\
	\vspace{0.1cm}
	\fecha\\
	\vspace{0.4cm}
\end{center}

\begin{preguntas}
\item Determine el valor de $r$ de modo que $y(x) = e^{rx}$ sea solución de la ecuación:
	$$3y'' + 3y' - 4y=0$$
\begin{solucion}
En primer lugar, calculamos las primeras dos derivadas de $y(x)$
		$$y(x) = e^{rx}$$
		$$y'(x) = re^{rx}$$
		$$y''(x) = r^2e^{rx}$$
		y procedemos a reemplazar en la ecuación
		$$3r^2e^{rx} + 3re^{rx} - 4e^{rx} = 0$$
		$$3r^2 + 3r - 4 = 0$$
		$$r = \dfrac{-3 \pm \sqrt[]{57}}{6}$$
\end{solucion}
\item Calcule $\lim\limits_{x \ra 0+} x^{x}
\begin{solucion}
Para calcular este límite, se usa la propiedad siguiente
		$$e^{ln(algo)} = algo$$
		Utilizando esta propiedad, 
		$$\lim\limits_{x \ra 0+} x^{x} = \lim\limits_{x \ra 0+} e^{ln(x^x)} = \lim\limits_{x \ra 0+} e^{xln(x)}$$
		Como $e$ es una constante, tenemos que
		$$\lim\limits_{x \ra 0+} e^{xln(x)} = e^{\lim\limits_{x \ra 0+} xln(x)}$$
		Por lo que en primer lugar calcularemos el límite del exponente, es decir
		$$\lim\limits_{x \ra 0+} xln(x) = \lim\limits_{x \ra 0+} \dfrac{ln(x)}{\frac{1}{x}} \stackrel{L'H}{=} \lim\limits_{x \ra 0+} \dfrac{\frac{1}{x}}{-\frac{1}{x^2}} = \lim\limits_{x \ra 0+} -x = 0 $$
\end{solucion}
\item Sea $a_n = \dfrac{n!}{n^n}$, calcular $\lim\limits_{n \ra \infty} \dfrac{a_{n+1}}{a_n}$
\begin{solucion}
$$\lim\limits_{n \ra \infty} \dfrac{a_{n+1}}{a_n} = \lim\limits_{n \ra \infty} \dfrac{\dfrac{(n+1)!}{(n+1)^{n+1}}}{\dfrac{n!}{n^n}} = \lim\limits_{n \ra \infty} \dfrac{(n+1)!}{n!} \cdot \dfrac{n^n}{(n+1)^{n+1}} =\lim\limits_{n \ra \infty} (n+1) \cdot \dfrac{n^n}{(n+1)^{n+1}}$$
		$$=\lim\limits_{n \ra \infty} \dfrac{n^n}{(n+1)^{n}} = \lim\limits_{n \ra \infty} \left(\dfrac{n}{n+1}\right)^n = \lim\limits_{n \ra \infty} \dfrac{1}{\left(\dfrac{n+1}{n}\right)^n} = \lim\limits_{n \ra \infty} \dfrac{1}{\left(1+\dfrac{1}{n}\right)^n}$$
		$$ = \dfrac{1}{e}$$
\end{solucion}
\item Calcular $F'(0)$, siendo
	$$F(x) = \displaystyle\int_0^{5x+1} \dfrac{e^{t^2}}{1+t^4}dt$$
\begin{solucion}
Usando el TFC, tenemos que
		$$F'(x) = \dfrac{e^{(5x+1)^2}}{1+(5x+1)^4} \cdot 5$$
		Luego, evaluamos en $x=0$
		$$F'(0) = \dfrac{5e^{(1)^2}}{1+(1)^4} = \dfrac{5e}{2}$$
\end{solucion}
\item Resolver las siguientes integrales indefinidas
\begin{tasks}(2)
\task $\displaystyle\int (3x^2 + 2x + 1)dx$
\task $\displaystyle\int x(x+1)(x+2)dx$
\task $\displaystyle\int sen^2(x)dx$
\task $\displaystyle\int (1+e)^xdx$
\end{tasks}
\begin{solucion}

\begin{enumerate}[a)]
\item $\displaystyle\int (3x^2 + 2x + 1)dx = \displaystyle\int 3x^2dx + \displaystyle\int 2xdx + \displaystyle\int dx  = 3\displaystyle\int x^2dx + 2\displaystyle\int xdx + \displaystyle\int dx $
			$$= 3\dfrac{x^3}{3} + 2\dfrac{x^2}{2} + x + c = x^3 + x^2 + x + c$$
\item $\displaystyle\int x(x+1)(x+2)dx = \displaystyle\int (x^3+3x^2+2x)dx = \dfrac{x^4}{4}+x^3+x+c$
\item $\displaystyle\int sen^2(x)dx = \displaystyle\int \dfrac{1-cos(2x)}{2}dx = \displaystyle\int \dfrac{1}{2}dx - \displaystyle\int \dfrac{cos(2x)}{2} = \dfrac{x}{2} -\dfrac{sen(2x)}{4} + c$
\item $\displaystyle\int (1+e)^xdx = \dfrac{(1+e)^x}{ln(1+e)}$
\end{enumerate}
\end{solucion}
\item Resolver las siguientes integrales indefinidas
\begin{tasks}(2)
\task $\displaystyle\int \dfrac{e^{ln(x)}}{x^2+7}dx$
\task $\displaystyle\int (x+2)sen(x^2+4x-6)dx$
\task $\displaystyle\int \dfrac{3x}{\sqrt[3]{x^2+3}}dx$
\task $\displaystyle\int \dfrac{e^x-1}{e^x+1}dx$
\end{tasks}
\begin{solucion}

\begin{enumerate}[a)]
\item $\displaystyle\int \dfrac{e^{ln(x)}}{x^2+7}dx = \displaystyle\int \dfrac{x}{x^2+7}dx$\\
			Usamos la sustitución $u = x^2+7 \ra du = 2xdx$
			$$= \dfrac{1}{2} \displaystyle\int \dfrac{2}{u}dx = \dfrac{1}{2} ln(u) + c = \dfrac{1}{2}ln(x^2+7) + c$$
\item $\displaystyle\int (x+2)sen(x^2+4x-6)dx$\\
			Utilizando la sustitución $u = x^2+4x-6 \ra du = 2x + 4 \ra du = 2(x+2)$
			$$=\dfrac{1}{2}\displaystyle\int sen(u)du = -\dfrac{1}{2}cos(u) = -\dfrac{1}{2}cos(x^2+4x-6)+c$$
\item $\displaystyle\int \dfrac{3x}{\sqrt[3]{x^2+3}}dx$\\
			La sustitución a usar es $u=x^2+3 \ra du=2xdx$
			$$= \dfrac{3}{2} \displaystyle\int \dfrac{1}{\sqrt[3]{u}}du = \dfrac{3}{2} \displaystyle\int \dfrac{1}{u^{1/3}}du = \dfrac{3}{2} \displaystyle\int u^{-1/3}du = \dfrac{3}{2}\dfrac{u^{2/3}}{\frac{2}{3}} + c = \dfrac{9}{4}(t^2+3)^{2/3}+c$$
\item $\displaystyle\int \dfrac{e^x-1}{e^x+1}dx = \displaystyle\int \dfrac{e^x}{e^x+1}dx - \displaystyle\int \dfrac{1}{e^x+1}dx$\\
			Resolvamos ambas por separado. En primer lugar,
			$$ \displaystyle\int \dfrac{e^x}{e^x+1}dx$$
			Utilizamos la sustitución $u=e^x+1 \ra du=e^x$
			$$= \displaystyle\int \dfrac{du}{u} = ln(u) + c = ln(e^x+1)+c$$
			Ahora, resolveremos
			$$-\displaystyle\int \dfrac{1}{e^x+1}dx = -\displaystyle\int \dfrac{e^xe^{-x}}{e^x+1}dx= -\displaystyle\int \dfrac{e^{-x}}{e^{-x}(e^x+1)}dx = -\displaystyle\int \dfrac{e^{-x}}{1+e^{-x}}dx$$
			Usando $u=1+e^{-x} \ra du = -e^{-x}$
			$$= \displaystyle\int \dfrac{du}{u}dx = ln(u) +c = ln(1+e^{-x})+c$$
			Por último, al sumar ambas obtenemos
			$$\displaystyle\int \dfrac{e^x-1}{e^x+1}dx = ln(1+e^x) + ln(1+e^{-x}) + c = ln((1+e^x)(1+e^{-x}))+c$$
\end{enumerate}
\end{solucion}
\item Resolver las siguientes integrales definidas
\begin{tasks}(2)
\task $\displaystyle\int_0^2 \dfrac{x^2}{x^3+8} dx$
\task $\displaystyle\int_{-1}^1 xsen(1-x^2)dx$
\end{tasks}
\begin{solucion}

\begin{enumerate}[a)]
\item $\displaystyle\int_0^2 \dfrac{x^2}{x^3+8} dx$\\
			Usando la sustitución
			$$u=x^3+8 \ra du =  3x^2dx$$
			$$x \in (0, 2) \ra u \in (8, 16)$$
			$$\displaystyle\int_0^2 \dfrac{x^2}{x^3+8}dx = \dfrac{1}{3}\displaystyle\int_8^{16} \dfrac{du}{u} = \dfrac{1}{3}ln(u) \ev_8^{16} = \dfrac{1}{3}(ln(16)-ln(8)) = \dfrac{1}{3}ln\left(\dfrac{16}{8}\right) = \dfrac{1}{3}ln(2)$$
\item $\displaystyle\int_{-1}^1 xsen(1-x^2)dx$\\
			Usamos el cambio de variable
			$$u = 1-x^2 \ra du = -2xdx$$
			$$x \in (-1, 1) \ra u \in (0, 0)$$
			Como la variable $u$ va de 0 a 0, entonces la integral es igual a 0. Luego,
			$$\displaystyle\int_{-1}^1 xsen(1-x^2)dx = 0$$
\end{enumerate}
\end{solucion}
\end{preguntas}
\end{document}